\chapter[Lorentz decomposition]{Irreducible Lorentz Decomposition of Fluctuations}
\label{ch:lorentz-decomposition}

\paragraph{Learning outcomes.}
After this chapter the reader should be able to:
\begin{enumerate}
  \item convert a fluctuation \(\delta\Phi^I{}_A\) to a spacetime
        tensor using the background frame;
  \item split that tensor into trace, antisymmetric, and symmetric
        trace-free pieces, and count components;
  \item match the split to the schematic formula used in the source
        notes,
        \(\delta\Phi^I{}_A=\tfrac14\eta_{\mu\nu}\delta g^{\mu\nu}I
        +A_\mu\gamma^\mu+T_{\mu\nu}\);
  \item distinguish algebraic irreducibility from the later
        dynamical conditions \(\nabla^\mu T_{\mu\nu}=0\) and
        \(m_T^2=0\);
  \item enter Chapter~6 knowing which piece is the candidate graviton.
\end{enumerate}

Chapter~4 produced a condensed background frame
\(\bar e^A_\mu=\delta^A_\mu\) and a composite Minkowski metric.
This chapter is linear algebra on the tangent space. No entanglement
condition is used. No mass gap is shown to vanish. The only task is
to name the pieces of \(\delta\Phi\).

\section{Fluctuations about the polar vacuum}
\label{sec:ch5-fluct}

Start from the polar field of Definition~\ref{def:vierbein},
\begin{equation}
  \label{eq:polar-again}
  \Phi^{IA}(x)=\bigl(v+\sigma(x)\bigr)\,\Lambda^A{}_I(x),
\end{equation}
and expand about the Minkowski condensate
\begin{equation}
  \label{eq:polar-bg}
  \langle\sigma\rangle=0,
  \qquad
  \langle\Lambda^A{}_I\rangle=\delta^A{}_I,
  \qquad
  \bar e^A_\mu=\delta^A_\mu.
\end{equation}
Write
\begin{equation}
  \label{eq:delta-Phi}
  \Phi^{IA}
  =v\,\delta^{IA}
   +\delta\Phi^{IA},
\end{equation}
with \(\delta\Phi\) first order in \(\sigma\) and in the deviation of
\(\Lambda\) from the identity. For the rest of the linearized theory
it is convenient to lower the internal index with the background
identification \(I\leftrightarrow A\) and convert the remaining tangent
index to a coordinate index with \(\bar e\):
\begin{equation}
  \label{eq:convert}
  \delta\Phi_{\mu\nu}
  \equiv
  \bar e^A{}_\mu\,\delta\Phi_{IA}\,
  \bar e^I{}_\nu.
\end{equation}
When \(N=4\) and the vacuum of Assumption~\ref{ass:frame-lock} locks
every internal direction to a tangent direction, \(\delta\Phi_{\mu\nu}\)
is a genuine \(4\times 4\) spacetime tensor. If \(N>4\),
\eqref{eq:convert} captures only the locked \(4\times 4\) block. The
orthogonal internal block is a collection of extra Higgs fields. They
are assumed heavy and integrated out; they are not candidate gravitons.

\begin{assumption}[Locked block dominates the infrared]
  \label{ass:locked-block}
  Fluctuations in internal directions orthogonal to the four locked
  axes have masses of order \(v\) or \(\LUV\) and may be discarded
  from the infrared Lorentz decomposition.
\end{assumption}

\section{Algebraic split under the Lorentz group}
\label{sec:ch5-algebra}

At a fixed event, \(\delta\Phi_{\mu\nu}\) is an element of
\(\mathbb{R}^4\otimes\mathbb{R}^4\). The Lorentz group
\(\mathrm{SO}(3,1)\) acts on each factor. The tensor product
decomposes into irreducible pieces
\begin{equation}
  \label{eq:tensor-product}
  \mathbf{4}\otimes\mathbf{4}
  =\mathbf{1}\oplus\mathbf{6}\oplus\mathbf{9},
\end{equation}
namely a scalar, an antisymmetric tensor, and a symmetric traceless
tensor. Explicitly,
\begin{equation}
  \label{eq:algebraic-split}
  \delta\Phi_{\mu\nu}
  =\frac14\eta_{\mu\nu}\psi
   +F_{\mu\nu}
   +T_{\mu\nu},
\end{equation}
where
\begin{equation}
  \label{eq:pieces-def}
  \begin{aligned}
    \psi
    &=\eta^{\alpha\beta}\delta\Phi_{\alpha\beta},\\
    F_{\mu\nu}
    &=\frac12\bigl(\delta\Phi_{\mu\nu}-\delta\Phi_{\nu\mu}\bigr),\\
    T_{\mu\nu}
    &=\frac12\bigl(\delta\Phi_{\mu\nu}+\delta\Phi_{\nu\mu}\bigr)
      -\frac14\eta_{\mu\nu}\psi.
  \end{aligned}
\end{equation}
By construction
\begin{equation}
  \label{eq:STF}
  T_{\mu\nu}=T_{\nu\mu},
  \qquad
  \eta^{\mu\nu}T_{\mu\nu}=0,
  \qquad
  F_{\mu\nu}=-F_{\nu\mu}.
\end{equation}

Component count in four dimensions:
\begin{center}
\begin{tabular}{@{}lcc@{}}
\toprule
Piece & Symbol & Components \\
\midrule
trace scalar & \(\psi\) & 1 \\
antisymmetric tensor & \(F_{\mu\nu}\) & 6 \\
symmetric traceless tensor & \(T_{\mu\nu}\) & 9 \\
total & \(\delta\Phi_{\mu\nu}\) & 16 \\
\bottomrule
\end{tabular}
\end{center}
This is algebra. It uses no field equation and no gauge condition.

\section{Translation into the notation of the source notes}
\label{sec:ch5-dictionary}

The microscopic notes write
\begin{equation}
  \label{eq:notes-split}
  \delta\Phi^I{}_A
  =\frac14\eta_{\mu\nu}\delta g^{\mu\nu}\,I
   +A_\mu\gamma^\mu
   +T_{\mu\nu}.
\end{equation}
Equation~\eqref{eq:notes-split} is~\eqref{eq:algebraic-split} after
three identifications.

\paragraph{Trace piece.}
On the linearized background of Section~\ref{sec:ch4-minkowski},
\begin{equation}
  \label{eq:psi-h}
  \psi=\eta^{\mu\nu}h_{\mu\nu}+O(\omega)
  =\delta g^{\mu}{}_{\mu}.
\end{equation}
The first term of~\eqref{eq:notes-split} is therefore the conformal
(trace) mode of the composite metric. The identity matrix \(I\) in
the notes is the locked vacuum \(\delta^I{}_A\), converted to
\(\eta_{\mu\nu}\) by~\eqref{eq:convert}.

\paragraph{Antisymmetric / vector piece.}
An antisymmetric tensor \(F_{\mu\nu}\) in four dimensions is
equivalent, via the Clifford algebra generated by matrices
\(\gamma^\mu\) obeying
\begin{equation}
  \label{eq:clifford}
  \{\gamma^\mu,\gamma^\nu\}=2\eta^{\mu\nu},
\end{equation}
to a vector-plus-axial combination. The notes keep only the vector
current \(A_\mu\gamma^\mu\), which is the part that couples to the
internal gauge field of~\eqref{eq:Dmu-Phi}. The remaining three
(axial / magnetic) components of \(F_{\mu\nu}\) are local Lorentz
parameters. They belong with the Goldstone \(\omega_{AB}\) of
\eqref{eq:linear-frame}, not with a new propagating vector.

This book uses the following official translation:
\begin{equation}
  \label{eq:F-split}
  F_{\mu\nu}
  =\nabla_\mu A_\nu-\nabla_\nu A_\mu
   +\omega_{\mu\nu},
\end{equation}
with \(A_\mu\) the internal gauge current and \(\omega_{\mu\nu}\) the
local Lorentz Goldstone. Chapter~10 eats \(\omega\). The Maxwell-like
field strength of \(A_\mu\) is a standard gauge excitation and is not
a graviton.

\paragraph{Symmetric traceless piece.}
The tensor \(T_{\mu\nu}\) in~\eqref{eq:notes-split} is identical to
the \(T_{\mu\nu}\) of~\eqref{eq:pieces-def}. It is the only algebraic
piece with the index symmetries of a graviton.

\section{What \(T_{\mu\nu}\) is not yet}
\label{sec:ch5-not-yet}

The source notes also write
\begin{equation}
  \label{eq:notes-extra}
  T^\mu{}_\mu=0,
  \qquad
  \nabla^\mu T_{\mu\nu}=0.
\end{equation}
The first of these is~\eqref{eq:STF} and has already been imposed.
The second is not algebraic. It is four differential conditions.

Those four conditions are the linearized statement that \(T_{\mu\nu}\)
is taken in de Donder / Lorentz-like gauge, or, on shell, that a
massive spin-2 field has five rather than nine components. They are
not free to impose until one knows the kinetic operator of \(T_{\mu\nu}\).

\begin{center}
\begin{tabular}{@{}lll@{}}
\toprule
Statement & Type & Status in this chapter \\
\midrule
\(T_{\mu\nu}=T_{\nu\mu}\) & algebraic & imposed \\
\(T^\mu{}_\mu=0\) & algebraic & imposed \\
\(\nabla^\mu T_{\mu\nu}=0\) & gauge or on-shell & \emph{not} imposed \\
\(m_T^2=0\) & dynamical & Chapter~7 \\
\(\square T_{\mu\nu}=0\) with two TT modes & dynamical & Chapter~7 \\
\bottomrule
\end{tabular}
\end{center}

A massive symmetric traceless tensor in four dimensions has five
on-shell degrees of freedom (the Fierz--Pauli count). A massless one
with the gauge invariance
\begin{equation}
  \label{eq:tt-gauge}
  \delta_\xi T_{\mu\nu}
  =\nabla_\mu\xi_\nu+\nabla_\nu\xi_\mu
   -\tfrac12\eta_{\mu\nu}\nabla\cdot\xi
\end{equation}
has two. The difference is the entire physical question of
Chapters~6 and~7. It cannot be settled by rearranging indices.

\section{Kinetic terms before any entanglement condition}
\label{sec:ch5-kinetic}

Inserting~\eqref{eq:algebraic-split} into the quadratic expansion of
\(S_{\mathrm{micro}}\) produces, schematically,
\begin{equation}
  \label{eq:S2-split}
  S^{(2)}
  =S^{(2)}_\psi[\psi]
   +S^{(2)}_A[A]
   +S^{(2)}_\omega[\omega]
   +S^{(2)}_T[T]
   +S^{(2)}_{\mathrm{mix}}.
\end{equation}
On a Minkowski condensate the mixing terms that survive after gauge
fixing are of the form \(\psi\,\partial\partial T\) and
\(A\,\partial T\). They are the usual linearized scalar--tensor and
vector--tensor mixings. Two structural facts will be used later.

\begin{enumerate}
  \item The tensor block has the generic form
  \begin{equation}
    \label{eq:ST-generic}
    S^{(2)}_T
    =\int\mathrm{d}^4x
     \Bigl(
       \tfrac12\nabla_\rho T_{\mu\nu}\nabla^\rho T^{\mu\nu}
       -\tfrac12 m_T^2 T_{\mu\nu}T^{\mu\nu}
       +\cdots
     \Bigr),
  \end{equation}
  where the ellipsis contains the additional derivative structures
  required by Fierz--Pauli if a mass term is present (the
  \(T^\mu{}_\mu\) piece has already been projected out). The mass
  \(m_T^2\) is a function of \(v\), \(\lambda\), and the condensate
  alignment. Nothing in Chapter~3 forces it to vanish.
  \item The conformal block \(S^{(2)}_\psi\) is the scalar-tensor
        sector. If \(\psi\) remains light, the infrared theory is not
        pure general relativity even if \(m_T=0\). Killing or gapping
        \(\psi\) is a separate requirement, recorded as such in
        Chapter~7 and again in Chapter~17.
\end{enumerate}

Equation~\eqref{eq:ST-generic} is the action whose mass term
Chapter~7 must kill. Writing it here prevents that later argument
from looking as if \(T_{\mu\nu}\) arrived already massless.

\section{Diffeomorphism action on the pieces}
\label{sec:ch5-diffs}

An infinitesimal diffeomorphism generated by \(\xi^\mu\) acts on the
composite metric as
\begin{equation}
  \label{eq:lie-g}
  \delta_\xi g_{\mu\nu}
  =\nabla_\mu\xi_\nu+\nabla_\nu\xi_\mu.
\end{equation}
On the linearized Minkowski background this is
\(\delta_\xi h_{\mu\nu}=\partial_\mu\xi_\nu+\partial_\nu\xi_\mu\).
Passing through~\eqref{eq:algebraic-split},
\begin{equation}
  \label{eq:lie-pieces}
  \begin{aligned}
    \delta_\xi\psi
    &=2\partial\cdot\xi,\\
    \delta_\xi F_{\mu\nu}
    &=0
     \quad\text{(to first order, after the symmetric projection)},\\
    \delta_\xi T_{\mu\nu}
    &=\partial_\mu\xi_\nu+\partial_\nu\xi_\mu
      -\tfrac12\eta_{\mu\nu}\partial\cdot\xi.
  \end{aligned}
\end{equation}
The last line is the gauge transformation that a massless spin-2
field must possess. If \(m_T^2\neq 0\), \eqref{eq:lie-pieces} is not
a symmetry of~\eqref{eq:ST-generic}, and the Ward identity that
Chapter~6 will write on a causal diamond is obstructed. That
obstruction is the lever used to force \(m_T\to 0\).

\section{Worked count: from sixteen components to two polarizations}
\label{sec:ch5-count}

The following table is the map the rest of the book will follow. Only
the first two columns are established in this chapter.

\begin{center}
\begin{tabular}{@{}llll@{}}
\toprule
Stage & What is removed & Remaining & Where \\
\midrule
algebraic split & --- & 16 & this chapter \\
eat \(\omega_{AB}\) & 6 Lorentz Goldstones & 10 & Ch.~10 \\
gap or constrain \(\psi\) & 1 conformal scalar & 9 & Ch.~7, 17 \\
gap \(A_\mu\) or eat it & 4 vector & 5 & Ch.~10 \\
impose \(\nabla^\mu T_{\mu\nu}=0\) & 4 constraints & 5 & on-shell FP \\
send \(m_T\to 0\) and gauge & 3 more & 2 TT & Ch.~6--7 \\
\bottomrule
\end{tabular}
\end{center}

The last line is the target advertised in Chapter~1. It is not a
result of the decomposition.

\section{Exercises}
\label{sec:ch5-exercises}

\begin{enumerate}
  \item \textbf{Projectors.}
    Write the three projectors on \(\mathbb{R}^4\otimes\mathbb{R}^4\)
    that extract \(\psi\), \(F_{\mu\nu}\), and \(T_{\mu\nu}\) from
    \(\delta\Phi_{\mu\nu}\). Verify they are complementary and
    idempotent.
  \item \textbf{Component count.}
    Confirm the \(1+6+9=16\) split by writing every independent
    component of a general real \(4\times 4\) matrix in an explicit
    basis.
  \item \textbf{Fierz--Pauli count.}
    A massive spin-2 field in \(d=4\) has how many on-shell degrees
    of freedom? Derive the number from \(\mathrm{SO}(3)\) little-group
    representations. Where does the number five sit in the table of
    Section~\ref{sec:ch5-count}?
  \item \textbf{Massless count.}
    Repeat for a massless spin-2 field with gauge invariance
    \eqref{eq:tt-gauge}. Recover two helicities.
  \item \textbf{Trace versus conformal mode.}
    From~\eqref{eq:linear-metric} and~\eqref{eq:psi-h}, show that
    \(\psi\) is the linearized conformal factor of \(g_{\mu\nu}\)
    plus a term involving \(\omega\). In what gauge does the extra
    term drop?
  \item \textbf{Diffeomorphism on \(T_{\mu\nu}\).}
    Derive the third line of~\eqref{eq:lie-pieces} from
    \eqref{eq:lie-g} and~\eqref{eq:pieces-def}.
  \item \textbf{\(N>4\).}
    Suppose \(N=8\). How many extra internal scalars appear outside
    the locked \(4\times 4\) block? Why are they not gravitons?
  \item \textbf{Danger of imposing \(\nabla^\mu T_{\mu\nu}=0\) early.}
    Take~\eqref{eq:ST-generic} with \(m_T\neq 0\) and impose
    \(\partial^\mu T_{\mu\nu}=0\) off shell as a constraint. Show
    that the constraint is consistent with the massive equation of
    motion, but that the massless limit \(m_T\to 0\) must be taken
    \emph{before} residual gauge transformations are used to reach
    two polarizations. This is why Chapter~7 cannot be skipped.
\end{enumerate}

\section{Bridge to Chapter~6}
\label{sec:ch5-bridge}

The candidate graviton has a name, \(T_{\mu\nu}\), nine algebraic
components, and a generic mass \(m_T\). The conformal and vector
pieces have also been named, and will be gapped or eaten rather than
identified with gravitational waves.

Nothing in the microscopic potential~\eqref{eq:V} sets \(m_T=0\).
The mechanism proposed by NAMCG is not a further tuning of
\(\lambda\). It is the entanglement-stationarity condition
\(\delta S_D=0\) on causal diamonds, imposed on the same fluctuations
just decomposed. Chapter~6 defines that condition from the modular
Hamiltonian of the matrix condensate and shows how it acts as a Ward
identity. Chapter~7 applies the identity to \(T_{\mu\nu}\) and
obtains \(m_T\to 0\).